% ============================================================
% fig_encoder_block.tex: Bloque Encoder detallado
% Uso: % ============================================================
% fig_encoder_block.tex: Bloque Encoder detallado
% Uso: % ============================================================
% fig_encoder_block.tex: Bloque Encoder detallado
% Uso: % ============================================================
% fig_encoder_block.tex: Bloque Encoder detallado
% Uso: \input{figures/fig_encoder_block}
% ============================================================
\begin{figure}[H]
\centering
\begin{tikzpicture}[
    block/.style={draw, rounded corners, fill=blue!8, draw=blue!40,
                  minimum width=4.2cm, minimum height=0.7cm,
                  font=\small, align=center},
    norm/.style={draw, rounded corners, fill=orange!10, draw=orange!50,
                 minimum width=4.2cm, minimum height=0.55cm,
                 font=\small, align=center},
    add/.style={draw, circle, fill=yellow!15, draw=yellow!60!black,
                minimum size=0.5cm, font=\small\bfseries},
    arrow/.style={-{Stealth}, thick},
    skip/.style={-{Stealth}, thick, gray!60, dashed}
]

% ── Entrada ──────────────────────────────────────────────────
\node[font=\footnotesize, gray] (inp) at (0, 0) {$\bm{X}$ (embeddings + PE)};

% ── Sub-bloque 1: Multi-Head Self-Attention ──────────────────
\node[block] (mha)  at (0, 1.4) {Multi-Head\\Self-Attention};
\node[add]   (add1) at (0, 2.6) {$+$};
\node[norm]  (ln1)  at (0, 3.4) {Layer Norm};

% ── Sub-bloque 2: Feed-Forward ───────────────────────────────
\node[block] (ffn)  at (0, 4.6) {Feed-Forward\\Network};
\node[add]   (add2) at (0, 5.8) {$+$};
\node[norm]  (ln2)  at (0, 6.6) {Layer Norm};

% ── Salida ───────────────────────────────────────────────────
\node[font=\footnotesize, gray] (out) at (0, 7.6) {$\bm{Z}$ (representaciones contextualizadas)};

% ── Flechas principales ──────────────────────────────────────
\draw[arrow] (inp)  -- (mha);
\draw[arrow] (mha)  -- (add1);
\draw[arrow] (add1) -- (ln1);
\draw[arrow] (ln1)  -- (ffn);
\draw[arrow] (ffn)  -- (add2);
\draw[arrow] (add2) -- (ln2);
\draw[arrow] (ln2)  -- (out);

% ── Conexiones residuales ─────────────────────────────────────
% Residual 1: entrada → add1
\draw[skip] (inp.east) -- ++(1.6,0) -- ++(0,2.6) -- (add1.east);
\node[font=\tiny, gray] at (2.0, 1.3) {residual};

% Residual 2: después de ln1 → add2
\draw[skip] (ln1.east) -- ++(1.4,0) -- ++(0,2.3) -- (add2.east);

% ── Anotaciones laterales ─────────────────────────────────────
\node[draw=blue!20, fill=blue!4, rounded corners, font=\tiny,
      inner sep=4pt, align=left] at (-3.2, 1.4)
    {Cada token consulta\\a todos los dem\'as.\\Produce representaci\'on\\contextualizada.};
\draw[-{Stealth}, gray!50, thin] (-1.8, 1.4) -- (-2.1, 1.4);

\node[draw=green!20, fill=green!4, rounded corners, font=\tiny,
      inner sep=4pt, align=left] at (-3.2, 4.6)
    {Transf. no lineal\\aplicada a cada token\\independientemente.\\ReLU o GELU.};
\draw[-{Stealth}, gray!50, thin] (-1.8, 4.6) -- (-2.1, 4.6);

\node[draw=orange!20, fill=orange!4, rounded corners, font=\tiny,
      inner sep=4pt, align=left] at (3.2, 3.4)
    {Normaliza activaciones.\\Estabiliza entrenamiento.};
\draw[-{Stealth}, gray!50, thin] (1.8, 3.4) -- (2.1, 3.4);

% ── Corchete del bloque ───────────────────────────────────────
\draw[gray!50, thick]
    (-2.6, 1.0) -- (-2.8, 1.0) -- (-2.8, 7.0) -- (-2.6, 7.0);
\node[font=\footnotesize, gray, rotate=90] at (-3.1, 4.0)
    {$\times N$ bloques};

\end{tikzpicture}
\caption{Estructura de un bloque encoder. Cada bloque contiene dos
sub-capas: atenci\'on multi-cabeza y una red feed-forward. Ambas est\'an
rodeadas por una conexi\'on residual (\textit{skip connection}) y una
normalizaci\'on de capa (\textit{Layer Norm}). El bloque se apila $N$ veces
para aumentar la capacidad representacional del modelo~\cite{vaswani2017attention}.}
\label{fig:encoder_block}
\end{figure}

% ============================================================
\begin{figure}[H]
\centering
\begin{tikzpicture}[
    block/.style={draw, rounded corners, fill=blue!8, draw=blue!40,
                  minimum width=4.2cm, minimum height=0.7cm,
                  font=\small, align=center},
    norm/.style={draw, rounded corners, fill=orange!10, draw=orange!50,
                 minimum width=4.2cm, minimum height=0.55cm,
                 font=\small, align=center},
    add/.style={draw, circle, fill=yellow!15, draw=yellow!60!black,
                minimum size=0.5cm, font=\small\bfseries},
    arrow/.style={-{Stealth}, thick},
    skip/.style={-{Stealth}, thick, gray!60, dashed}
]

% ── Entrada ──────────────────────────────────────────────────
\node[font=\footnotesize, gray] (inp) at (0, 0) {$\bm{X}$ (embeddings + PE)};

% ── Sub-bloque 1: Multi-Head Self-Attention ──────────────────
\node[block] (mha)  at (0, 1.4) {Multi-Head\\Self-Attention};
\node[add]   (add1) at (0, 2.6) {$+$};
\node[norm]  (ln1)  at (0, 3.4) {Layer Norm};

% ── Sub-bloque 2: Feed-Forward ───────────────────────────────
\node[block] (ffn)  at (0, 4.6) {Feed-Forward\\Network};
\node[add]   (add2) at (0, 5.8) {$+$};
\node[norm]  (ln2)  at (0, 6.6) {Layer Norm};

% ── Salida ───────────────────────────────────────────────────
\node[font=\footnotesize, gray] (out) at (0, 7.6) {$\bm{Z}$ (representaciones contextualizadas)};

% ── Flechas principales ──────────────────────────────────────
\draw[arrow] (inp)  -- (mha);
\draw[arrow] (mha)  -- (add1);
\draw[arrow] (add1) -- (ln1);
\draw[arrow] (ln1)  -- (ffn);
\draw[arrow] (ffn)  -- (add2);
\draw[arrow] (add2) -- (ln2);
\draw[arrow] (ln2)  -- (out);

% ── Conexiones residuales ─────────────────────────────────────
% Residual 1: entrada → add1
\draw[skip] (inp.east) -- ++(1.6,0) -- ++(0,2.6) -- (add1.east);
\node[font=\tiny, gray] at (2.0, 1.3) {residual};

% Residual 2: después de ln1 → add2
\draw[skip] (ln1.east) -- ++(1.4,0) -- ++(0,2.3) -- (add2.east);

% ── Anotaciones laterales ─────────────────────────────────────
\node[draw=blue!20, fill=blue!4, rounded corners, font=\tiny,
      inner sep=4pt, align=left] at (-3.2, 1.4)
    {Cada token consulta\\a todos los dem\'as.\\Produce representaci\'on\\contextualizada.};
\draw[-{Stealth}, gray!50, thin] (-1.8, 1.4) -- (-2.1, 1.4);

\node[draw=green!20, fill=green!4, rounded corners, font=\tiny,
      inner sep=4pt, align=left] at (-3.2, 4.6)
    {Transf. no lineal\\aplicada a cada token\\independientemente.\\ReLU o GELU.};
\draw[-{Stealth}, gray!50, thin] (-1.8, 4.6) -- (-2.1, 4.6);

\node[draw=orange!20, fill=orange!4, rounded corners, font=\tiny,
      inner sep=4pt, align=left] at (3.2, 3.4)
    {Normaliza activaciones.\\Estabiliza entrenamiento.};
\draw[-{Stealth}, gray!50, thin] (1.8, 3.4) -- (2.1, 3.4);

% ── Corchete del bloque ───────────────────────────────────────
\draw[gray!50, thick]
    (-2.6, 1.0) -- (-2.8, 1.0) -- (-2.8, 7.0) -- (-2.6, 7.0);
\node[font=\footnotesize, gray, rotate=90] at (-3.1, 4.0)
    {$\times N$ bloques};

\end{tikzpicture}
\caption{Estructura de un bloque encoder. Cada bloque contiene dos
sub-capas: atenci\'on multi-cabeza y una red feed-forward. Ambas est\'an
rodeadas por una conexi\'on residual (\textit{skip connection}) y una
normalizaci\'on de capa (\textit{Layer Norm}). El bloque se apila $N$ veces
para aumentar la capacidad representacional del modelo~\cite{vaswani2017attention}.}
\label{fig:encoder_block}
\end{figure}

% ============================================================
\begin{figure}[H]
\centering
\begin{tikzpicture}[
    block/.style={draw, rounded corners, fill=blue!8, draw=blue!40,
                  minimum width=4.2cm, minimum height=0.7cm,
                  font=\small, align=center},
    norm/.style={draw, rounded corners, fill=orange!10, draw=orange!50,
                 minimum width=4.2cm, minimum height=0.55cm,
                 font=\small, align=center},
    add/.style={draw, circle, fill=yellow!15, draw=yellow!60!black,
                minimum size=0.5cm, font=\small\bfseries},
    arrow/.style={-{Stealth}, thick},
    skip/.style={-{Stealth}, thick, gray!60, dashed}
]

% ── Entrada ──────────────────────────────────────────────────
\node[font=\footnotesize, gray] (inp) at (0, 0) {$\bm{X}$ (embeddings + PE)};

% ── Sub-bloque 1: Multi-Head Self-Attention ──────────────────
\node[block] (mha)  at (0, 1.4) {Multi-Head\\Self-Attention};
\node[add]   (add1) at (0, 2.6) {$+$};
\node[norm]  (ln1)  at (0, 3.4) {Layer Norm};

% ── Sub-bloque 2: Feed-Forward ───────────────────────────────
\node[block] (ffn)  at (0, 4.6) {Feed-Forward\\Network};
\node[add]   (add2) at (0, 5.8) {$+$};
\node[norm]  (ln2)  at (0, 6.6) {Layer Norm};

% ── Salida ───────────────────────────────────────────────────
\node[font=\footnotesize, gray] (out) at (0, 7.6) {$\bm{Z}$ (representaciones contextualizadas)};

% ── Flechas principales ──────────────────────────────────────
\draw[arrow] (inp)  -- (mha);
\draw[arrow] (mha)  -- (add1);
\draw[arrow] (add1) -- (ln1);
\draw[arrow] (ln1)  -- (ffn);
\draw[arrow] (ffn)  -- (add2);
\draw[arrow] (add2) -- (ln2);
\draw[arrow] (ln2)  -- (out);

% ── Conexiones residuales ─────────────────────────────────────
% Residual 1: entrada → add1
\draw[skip] (inp.east) -- ++(1.6,0) -- ++(0,2.6) -- (add1.east);
\node[font=\tiny, gray] at (2.0, 1.3) {residual};

% Residual 2: después de ln1 → add2
\draw[skip] (ln1.east) -- ++(1.4,0) -- ++(0,2.3) -- (add2.east);

% ── Anotaciones laterales ─────────────────────────────────────
\node[draw=blue!20, fill=blue!4, rounded corners, font=\tiny,
      inner sep=4pt, align=left] at (-3.2, 1.4)
    {Cada token consulta\\a todos los dem\'as.\\Produce representaci\'on\\contextualizada.};
\draw[-{Stealth}, gray!50, thin] (-1.8, 1.4) -- (-2.1, 1.4);

\node[draw=green!20, fill=green!4, rounded corners, font=\tiny,
      inner sep=4pt, align=left] at (-3.2, 4.6)
    {Transf. no lineal\\aplicada a cada token\\independientemente.\\ReLU o GELU.};
\draw[-{Stealth}, gray!50, thin] (-1.8, 4.6) -- (-2.1, 4.6);

\node[draw=orange!20, fill=orange!4, rounded corners, font=\tiny,
      inner sep=4pt, align=left] at (3.2, 3.4)
    {Normaliza activaciones.\\Estabiliza entrenamiento.};
\draw[-{Stealth}, gray!50, thin] (1.8, 3.4) -- (2.1, 3.4);

% ── Corchete del bloque ───────────────────────────────────────
\draw[gray!50, thick]
    (-2.6, 1.0) -- (-2.8, 1.0) -- (-2.8, 7.0) -- (-2.6, 7.0);
\node[font=\footnotesize, gray, rotate=90] at (-3.1, 4.0)
    {$\times N$ bloques};

\end{tikzpicture}
\caption{Estructura de un bloque encoder. Cada bloque contiene dos
sub-capas: atenci\'on multi-cabeza y una red feed-forward. Ambas est\'an
rodeadas por una conexi\'on residual (\textit{skip connection}) y una
normalizaci\'on de capa (\textit{Layer Norm}). El bloque se apila $N$ veces
para aumentar la capacidad representacional del modelo~\cite{vaswani2017attention}.}
\label{fig:encoder_block}
\end{figure}

% ============================================================
\begin{figure}[H]
\centering
\begin{tikzpicture}[
    block/.style={draw, rounded corners, fill=blue!8, draw=blue!40,
                  minimum width=4.2cm, minimum height=0.7cm,
                  font=\small, align=center},
    norm/.style={draw, rounded corners, fill=orange!10, draw=orange!50,
                 minimum width=4.2cm, minimum height=0.55cm,
                 font=\small, align=center},
    add/.style={draw, circle, fill=yellow!15, draw=yellow!60!black,
                minimum size=0.5cm, font=\small\bfseries},
    arrow/.style={-{Stealth}, thick},
    skip/.style={-{Stealth}, thick, gray!60, dashed}
]

% ── Entrada ──────────────────────────────────────────────────
\node[font=\footnotesize, gray] (inp) at (0, 0) {$\bm{X}$ (embeddings + PE)};

% ── Sub-bloque 1: Multi-Head Self-Attention ──────────────────
\node[block] (mha)  at (0, 1.4) {Multi-Head\\Self-Attention};
\node[add]   (add1) at (0, 2.6) {$+$};
\node[norm]  (ln1)  at (0, 3.4) {Layer Norm};

% ── Sub-bloque 2: Feed-Forward ───────────────────────────────
\node[block] (ffn)  at (0, 4.6) {Feed-Forward\\Network};
\node[add]   (add2) at (0, 5.8) {$+$};
\node[norm]  (ln2)  at (0, 6.6) {Layer Norm};

% ── Salida ───────────────────────────────────────────────────
\node[font=\footnotesize, gray] (out) at (0, 7.6) {$\bm{Z}$ (representaciones contextualizadas)};

% ── Flechas principales ──────────────────────────────────────
\draw[arrow] (inp)  -- (mha);
\draw[arrow] (mha)  -- (add1);
\draw[arrow] (add1) -- (ln1);
\draw[arrow] (ln1)  -- (ffn);
\draw[arrow] (ffn)  -- (add2);
\draw[arrow] (add2) -- (ln2);
\draw[arrow] (ln2)  -- (out);

% ── Conexiones residuales ─────────────────────────────────────
% Residual 1: entrada → add1
\draw[skip] (inp.east) -- ++(1.6,0) -- ++(0,2.6) -- (add1.east);
\node[font=\tiny, gray] at (2.0, 1.3) {residual};

% Residual 2: después de ln1 → add2
\draw[skip] (ln1.east) -- ++(1.4,0) -- ++(0,2.3) -- (add2.east);

% ── Anotaciones laterales ─────────────────────────────────────
\node[draw=blue!20, fill=blue!4, rounded corners, font=\tiny,
      inner sep=4pt, align=left] at (-3.2, 1.4)
    {Cada token consulta\\a todos los dem\'as.\\Produce representaci\'on\\contextualizada.};
\draw[-{Stealth}, gray!50, thin] (-1.8, 1.4) -- (-2.1, 1.4);

\node[draw=green!20, fill=green!4, rounded corners, font=\tiny,
      inner sep=4pt, align=left] at (-3.2, 4.6)
    {Transf. no lineal\\aplicada a cada token\\independientemente.\\ReLU o GELU.};
\draw[-{Stealth}, gray!50, thin] (-1.8, 4.6) -- (-2.1, 4.6);

\node[draw=orange!20, fill=orange!4, rounded corners, font=\tiny,
      inner sep=4pt, align=left] at (3.2, 3.4)
    {Normaliza activaciones.\\Estabiliza entrenamiento.};
\draw[-{Stealth}, gray!50, thin] (1.8, 3.4) -- (2.1, 3.4);

% ── Corchete del bloque ───────────────────────────────────────
\draw[gray!50, thick]
    (-2.6, 1.0) -- (-2.8, 1.0) -- (-2.8, 7.0) -- (-2.6, 7.0);
\node[font=\footnotesize, gray, rotate=90] at (-3.1, 4.0)
    {$\times N$ bloques};

\end{tikzpicture}
\caption{Estructura de un bloque encoder. Cada bloque contiene dos
sub-capas: atenci\'on multi-cabeza y una red feed-forward. Ambas est\'an
rodeadas por una conexi\'on residual (\textit{skip connection}) y una
normalizaci\'on de capa (\textit{Layer Norm}). El bloque se apila $N$ veces
para aumentar la capacidad representacional del modelo~\cite{vaswani2017attention}.}
\label{fig:encoder_block}
\end{figure}
